\documentclass[10pt]{article}
\usepackage[margin=20mm]{geometry}
\usepackage{amsmath,amssymb,fontspec,unicode-math,xurl,hyperref}
\IfFontExistsTF{DejaVu Serif}{\setmainfont{DejaVu Serif}}{\setmainfont{Latin Modern Roman}}
\setmathfont{latinmodern-math.otf}
\setlength{\parindent}{0pt}\setlength{\parskip}{6pt}\sloppy
\title{Initial Scientific Review of
Adversarial Inception Backdoor Attacks against Reinforcement Learning}
\author{AI-assisted scientific review}
\date{24 October 2026}
\begin{document}
\maketitle
Abstract: The printed value-preservation argument fails when poisoned states can follow poisoned states. Under the admissible reward-history range explicitly used in the appendix, and after granting normalization and domain repairs, an exact finite construction gives a benign-state value \(15/16\) instead of \(1\). The supporting lemma and its base case are false under the stated interpretation, invalidating a derivation used for the central stealth guarantee. This supports Major without asserting that the final optimality theorem itself has been refuted or that the empirical results fail. Two local findings concern a doubly weighted transition constraint and an unguarded rollout-boundary index in the inspected release. Scientific adjudication remains pending.\par
\section*{Paper identity and reviewed version}
Ethan Rathbun, Alina Oprea, and Christopher Amato, Adversarial Inception Backdoor Attacks against Reinforcement Learning, ICML 2025, PMLR 267, pp. 51273–51296 [paper]. Assignment: 0160 / I25-43529. The reviewed published PDF contains 24 pages, including Appendix A. Page references below use one-based PDF pages, which agree with the internal printed page numbers. The frozen manifest records publisher publication on 6 October 2025 and retrieval on 24 October 2026. The inspected released source is pinned to `384f29901f2af6378649dff4c998b0fc6212989f` [release]. That pin identifies the available release; correspondence to the historical experimental executable is not established.\par
\section*{Scope and necessary background}
The paper proposes bounded-reward backdoor attacks against reinforcement learning. The theoretical construction changes transition and reward functions, while Q-Incept approximates the attack through changes to collected training trajectories. The abstract, main §4.3, and conclusion present simultaneous attack-success and stealth guarantees as central contributions.\par
Independent calculations concentrated on the formal guarantees and relevant source branches. The independent check checked the decisive original equations, proof dependency, reward-domain assumptions, and source lines. Original PDF pages 2, 5, and 16 were rendered and actually viewed. Numerical curves were not digitized, experiments were not rerun, and author code was neither imported nor executed.\par
Let \(S\) be the benign state set, \(S_p\) a disjoint poisoned copy, and \(\delta:S\to S_p\) a bijection. The action set contains the target \(a^+\). A discount \(0<\gamma<1\), poisoning probability \(0<\beta<1\), and policy \(\pi\) define the values under comparison. Equation (9), p. 5, remaps \(a^+\) at a poisoned state to an action maximizing a benign-state action value; other actions are unchanged. Equation (10) assigns probability \(1-\beta\) to benign successors and \(\beta\) to poisoned successors, including when the current state is poisoned.\par
Equation reference: eq:tau. With \(L=\inf R\), \(U=\sup R\), the published reward perturbation is \[
\tau(a,r,\widehat r)=
\begin{cases}
\min\{r+(\widehat r-L)/\gamma,U\}, & a=a^+,\\
\max\{r-(U-\widehat r)/\gamma,L\}, & a\ne a^+.
\end{cases}\] The prose identifies \(\widehat r\) as the previous received reward. Equation (12) adds a correction \(\gamma\mathbb E[r-\tau]\) only to benign-to-poisoned transitions; poisoned-origin transitions receive \(\tau\). Equation (7), p. 4, requires \(\inf R'\ge L\) and \(\sup R'\le U\). This is a condition to check, not a premise that can be omitted.\par
\section*{Substantive findings}
Problem 1 (Consecutive poisoned states break the value-preservation derivation).  Assessment: Major. Lemma 2, main p. 6 and Appendix A.1.4, p. 16, asserts, in paraphrase, \[V_\pi^{M'}(s)\ge V_\pi^M(s)
\qquad\text{for every }s\in S,\ \pi\in\Pi.\] Its base case asserts equality when action remapping is replaced by the identity. The step leading to Appendix equation (52) substitutes a poisoned action value by its benign counterpart plus \(\tau-r\). That substitution does not account for the different correction rewards when the successor is itself poisoned.\par
\subsection*{Explicit interpretation and granted repairs}
Three choices are stated before evaluating the claim. First, use the normalized full kernel in equation (10), granting the local normalization repair in Problem Problem 2 (\#prob:normalization). Second, pull a poisoned successor back through \(\delta^{-1}\) when a benign reward function is called on it. The poisoned-origin row of equation (12) does not explicitly write that successor pullback.\par
Third, use the admissible previous-reward domain \(\widehat r\in[L,U]\) explicitly assumed in Appendix equation (15), p. 13. In the benign-to-poisoned correction, evaluate the future \(\tau\) using the incoming reward of that next poisoned step. This is the forward-consistent interpretation of the prose’s cancellation of the next perturbation. The unmodified-incoming and corrected-received-incoming readings agree in the finite construction below; its consistency equation is verified explicitly.\par
These qualifications matter. Equations (11)–(12) print real argument domains without restricting \(\widehat r\). Taken literally on all real inputs, [eq:tau] (\#eq:tau) is not bounded: at \(L=0,U=1,\gamma=1/2,r=1/2\), \[\tau(a^+,1/2,-1)=-3/2,\qquad
\tau(a_0,1/2,2)=5/2 .\] Thus the example below is not claimed to satisfy an unrestricted-real-domain version of equation (7). It meets the bound on all state/action/successor triples for every admissible prior reward in \([0,1]\), under the stated incoming-reward interpretation. Other interpretations of the unspecified dot arguments require a complete definition; no invariance under every possible repair is asserted.\par
\subsection*{Finite model and verification of assumptions}
Equation reference: eq:benignreward. Take \[S=\{\ell,b,h\},\quad S_p=\{p_\ell,p_b,p_h\},\quad
\delta(q)=p_q,\quad A=\{a^+,a_0\},\quad
\beta=\gamma=\tfrac12 .\] The sets are finite and disjoint, and \(\delta\) is bijective. Both benign actions self-loop: \(T(q,a,q)=1\). Define the complete benign reward function, including impossible cross-component transitions, by \[
R(q,a,q')=r(q'),\qquad
r(\ell)=0,\quad r(b)=\tfrac12,\quad r(h)=1 .\] Consequently \(L=0,U=1\) are attained on actual self-loops. A uniform benign initial distribution, if needed, gives both extrema positive probability. The off-support definition in [eq:benignreward] (\#eq:benignreward) is part of this final construction; it strengthens the earlier auxiliary candidate without changing any operative benign transition.\par
Let \(\pi\) always take \(a_0\), at both benign and poisoned states. This is a valid policy in the universal quantifier. The argmax branch of \(\phi\) is never used on policy, and the actual remapping equals the identity. Off policy, the two benign actions have the same transition and reward functions, so action selection does not change those components.\par
For each \(q\), equation (10) gives the closed two-state transition matrix \[P_q=\begin{pmatrix}1/2&1/2\\1/2&1/2\end{pmatrix}
\quad\text{on }(q,p_q).\] In particular \(p_q\to p_q\) has positive probability. No inspected assumption excludes consecutive poisoned states; deleting that branch would change the published kernel.\par
Equation reference: eq:entry. The benign-to-poisoned entry rewards are \[c_\ell=0,\qquad c_b=\tfrac34,\qquad c_h=1.\] Each satisfies both the nominal incoming calculation and the received-reward consistency equation \[
c_q=r(q)+\tfrac12\bigl[r(q)-\tau(a_0,r(q),c_q)\bigr].\] For \(b\), \(\tau(a_0,1/2,3/4)=0\), giving \(c_b=3/4\). For \(\ell\) and \(h\), the corresponding perturbations are \(0\) and \(1\).\par
The reward constraint holds over all \(6\cdot2\cdot6=72\) state/action/successor triples, not only positive-probability transitions. Benign-to-benign rewards are \(r(\text{destination})\). Benign-to-poisoned rewards are the corresponding \(c_q\), all in \([0,1]\). For poisoned origins and any \(\widehat r\in[0,1]\), the target branch of [eq:tau] (\#eq:tau) lies in \([r,1]\) and the other branch in \([0,r]\), including actions not taken by the fixed policy. These clipped affine functions are nondecreasing in \(\widehat r\), so checking their endpoints also bounds the full interval. An exact calculation enumerated all 72 triples at both endpoints and verified [eq:entry] (\#eq:entry). Actual output: . An independent challenge reaches the same result in .\par
\subsection*{Correct values and failed substitution}
Equation reference: eq:vb, eq:vp, eq:counter. Start the central component from a benign state with a compatible prior reward \(1/2\). Its operative reward histories remain in \(\{0,1/2,3/4\}\), and \(\tau(a_0,1/2,\widehat r)=0\) for every such history. Thus the following two-state Bellman calculation is valid even if the previous reward is explicitly included in the underlying process: \[(R'_{b b},R'_{b p_b},R'_{p_b b},R'_{p_b p_b})
=\left(\tfrac12,\tfrac34,0,0\right).\] Writing \(v=V_\pi^{M'}(b)\), \(w=V_\pi^{M'}(p_b)\), with the latter understood on those consistent operative histories, \[\begin{aligned}
v&=\tfrac12(\tfrac12+\tfrac12v)+
   \tfrac12(\tfrac34+\tfrac12w)
 =\tfrac58+\tfrac14(v+w),\\
w&=\tfrac12(\tfrac12v)+\tfrac12(\tfrac12w)
 =\tfrac14(v+w).\end{aligned}\] Subtracting gives \(v-w=5/8\), while [eq:vp] (\#eq:vp) gives \(v=3w\). Therefore \[
v=\tfrac{15}{16}<1=\frac{1/2}{1-1/2}=V_\pi^M(b),
\qquad w=\tfrac5{16}.\] Both Bellman residuals are exactly zero in rational arithmetic. The original value is \(1\); the asserted nondecrease is violated by \(1/16\). The base-case equality also fails because remapping is already the identity on policy.\par
In this one-action-on-policy component, the equation-(52) substitution would require \[Q_\pi^{M'}(p_b,a_0)
=Q_\pi^{M'}(b,a_0)+\tau-r
=\tfrac{15}{16}-\tfrac12=\tfrac7{16}.\] The correct left side is \(5/16\), a residual of \(-1/8\). This identifies the failed step rather than merely reporting a different numerical value. An independent calculation solved the exact system and checked all three reward components. Its original reachable-history calculation is supplemented, not retroactively replaced, by the full admissible-domain check above.\par
More generally, within a constant-reward component with a consistent constant poisoned reward \(t\), the stated correction pattern gives \[v-w=(r-t)(1+\beta\gamma),\qquad
v=\frac{r}{1-\gamma}
-\frac{\beta^2\gamma^2(r-t)}{1-\gamma}.\] The residual comes from positive-probability consecutive poisoning. The example uses \(r=1/2,t=0,\beta=\gamma=1/2\).\par
\subsection*{Dependency, counterevidence, and repair limits}
Appendix A.1.5, p. 18, explicitly invokes Lemmas 1 and 2 in the forward proof of Theorem 2, which claims preservation of benign optimal policies. A false supporting lemma and invalid value substitution invalidate that central derivation. This is the reason for Major. The result does not depend on the proof’s appendix location or on demonstrating a benchmark decrease.\par
No direct counterexample to the final optimality biconditional is claimed here. The construction is policy dependent, and a valid replacement theorem would require its own fully specified reward/history semantics and proof. Other theoretical statements and empirical successes may survive.\par
The practical procedure is relevant counterevidence. Algorithm 1, p. 6, and the inspected release correct the previous sampled transition regardless of whether it was poisoned; p. 7 also explains a stronger practical reward choice. They are not identical to the three-row theoretical reward definition. Consequently the formal counterexample is not evidence that the reported Q-Incept attack-success or benign-return measurements are false.\par
In the finite central example, adding \(1/4\) to the \(p_b\to p_b\) reward gives \(v=1,w=1/2\); this diagnostic restores the local cancellation. It is not a proposed globally valid bounded-reward repair. The checker records this finite result separately. A full repair must specify history dependence and all transition types, recheck the global bound, and prove the resulting optimality statement.\par
Problem 2 (The transition constraint weights the mixture twice).  Assessment: Minor. Main §3, p. 2, equation (1), states the constraint \[(1-\beta)T'(s'\mid a,s)+
\beta T'(\delta(s')\mid a,s)=T(s'\mid a,s).\] At the same time \(T'\) is the full transition kernel on \(S\cup S_p\), as made explicit in equation (10), p. 5.\par
For a fixed row, let \(P_B=\sum_{s'\in S}T'(s'\mid a,s)\) and \(P_P=\sum_{s'\in S}T'(\delta(s')\mid a,s)\). Disjointness and bijectivity give \(P_B+P_P=1\). Summing the printed constraint instead requires \[(1-\beta)P_B+\beta P_P=1.\] For \(0<\beta<1\), its left side is at most \(\max\{\beta,1-\beta\}<1\), a contradiction. With \(\beta=1/2\), equation (10)’s valid row is \((1/2,1/2)\), but equation (1)’s left side is \(1/2\) rather than \(1\). Original pages 2 and 5 were visually checked.\par
The intended local repair is clear from the prose and equation (10): \[T'(s'\mid a,s)+T'(\delta(s')\mid a,s)=T(s'\mid a,s).\] Interpreting the displayed \(T'\) terms as branch-conditional kernels would require different notation and would conflict with their actual equation-(10) values. This is a supported local specification inconsistency. The report grants its repair in Problem Problem 1 (\#prob:value); it does not use the inconsistency to claim that every practical experiment is impossible. Subsequent probability expansions must be synchronized with whichever convention is adopted.\par
Problem 3 (The release uses a future rollout reward as the predecessor at index zero).  Assessment: Minor; historical impact Needs Evidence. In the pinned release, `adversary/OuterLoop.py`, lines 85–100, samples indices from a categorical distribution over the supplied states. Lines 140–150 then access `rewards[index-1]` without excluding index zero. The caller in `ppo.py`, lines 241–248, passes chronological per-environment rollout slices, including `rewards[:, i]`.\par
The code path is conditional on attack activation, the caller’s attack-success condition, and the selection budget. Index zero is included in the finite categorical support. Finite scores give it positive probability; positive selected scores cause the target action to be assigned. For example, finite two-action Q rows \((1,0),(0,0),(0,0)\), with the first action selected, yield a positive centered score at the first step and finite clipped logits. Such values are an illustrative legal branch condition, not a measured network output.\par
Take a three-step chronological reward slice \((1/2,1/2,1/2)\), discount \(1/2\), and cached empirical bounds \(L=0,U=1\) already established by earlier observations. The cached-bound condition is necessary: a first batch containing only \(1/2\) would otherwise initialize both bounds to \(1/2\). Assume index zero is selected for the target branch. The inspected expression sets its current reward to \(1\) and changes \[\texttt{rewards[-1]}
=\max\{0,\tfrac12-\tfrac12(1-\tfrac12)\}
=\tfrac14.\] The modified entry is the last future step of the supplied rollout. It is not the preceding transition, which lies outside that slice. The exact independent branch arithmetic appears in . No author function or learning environment was executed.\par
A larger rollout may reduce the frequency of selecting a particular boundary index, but it does not change this branch’s temporal meaning. No historical logs establish its occurrence or effect on published scores. A correction must explicitly handle the actual previous transition and episode/rollout boundaries; merely indexing the current slice cannot recover that transition. The general repair and its empirical impact remain unverified.\par
\section*{Overall impact and unresolved claims}
The supported overall assessment is Major, arising from Problem Problem 1 (\#prob:value)’s core theoretical dependency. Problems Problem 2 (\#prob:normalization) and Problem 3 (\#prob:boundary) remain local Minor findings with separate evidence. The finite model establishes a failure of the value-preservation claim under explicitly stated, source-supported semantics; it does not establish that all variants of adversarial inception fail.\par
The theoretical and practical algorithms differ in ways that prevent transferring the formal defect directly to measured attack success or benign return. Historical commit provenance, raw rollout logs, exact configuration-to-table correspondence, and reruns would be required to assess empirical effects. The full original text and selected release were inspected, but no comprehensive later-version or external-erratum search was performed. No priority or first-discovery claim is made.\par
\section*{Conclusion and state}
Finding | Assessment | Repair verification | Adoption | Remaining work\par
Value preservation | Major | Finite cancellation diagnostic only | Not requested | Full bounded formulation and proof\par
Mixture constraint | Minor | Aggregate normalization checked | Unknown | Consistent notation and expansions\par
Rollout predecessor | Minor | No implementation repair | Unknown | Boundary handling and historical effect\par
All scientific decisions remain pending. No human review has been supplied; `human\_reviews` is empty.\par
9 E. Rathbun, A. Oprea, and C. Amato, Adversarial Inception Backdoor Attacks against Reinforcement Learning, Proceedings of the 42nd International Conference on Machine Learning, PMLR 267 (2025), 51273–51296. Published proceedings PDF, 24 pages. https://proceedings.mlr.press/v267/rathbun25a.html. EthanRath, Backdoors-In-RL, inspected staged source at commit `384f29901f2af6378649dff4c998b0fc6212989f`. https://github.com/EthanRath/Backdoors-In-RL. Consulted source files and content hashes are recorded in the accompanying evidence bindings. The pin is not asserted to identify the historical experiment.\par
\section*{Actual checks and reusable evidence}
All evidence paths in this appendix are relative to the delivered bundle root. The input PDF and manifest are and . All 21 manifest-listed files matched their staged hashes during source preparation verification. Available original bytes were reused without source modification or new download.\par
The independent check executed the following reviewer-owned checks using the existing Windows Python and exact `fractions.Fraction` arithmetic:\par
> python -X utf8 verification/0160/parent\_check.py
> python -X utf8 verification/0160/parent\_domain\_check.py\par
The first checks the normalized transition row, exact Bellman solutions in the low/central/high components, operative previous-reward closure, the equation-(52) residual, a finite correction diagnostic, and the release’s boundary arithmetic. The second independently checks the strengthened destination-based reward completion over every state/action/successor triple, uses monotonicity to cover the whole admissible reward interval, verifies entry consistency, and recomputes the central Bellman residuals. Both completed successfully. Neither imports author code.\par
The independent scientific auxiliary wrote and its JSON output. A separate native challenge produced , , and later , its JSON output, and . The independent check read the decisive original text, code, calculations, challenge, and domain revision; agreement is not a substitute for those source checks.\par
The early auxiliary candidate completed impossible benign cross-component rewards by the source state’s reward. That completion preserves reachable calculations but can violate a literal all-state-pair bound off support. It is retained as a historical working candidate. Equation [eq:benignreward] (\#eq:benignreward) and the new full-domain check supersede that particular assumption. This revision must not be attributed to the earlier check. An initial auxiliary check also failed because it looked for the manifest under the source directory; its path was corrected before the successful run. A parent metadata read was truncated before JSON parsing; it was reread completely. Neither failure changed source bytes or supplies scientific evidence.\par
\section*{Source correspondence and hashes}
Static reading covered the relevant `OuterLoop.py` initialization, selection, reward-bound updates and predecessor writes, and the `ppo.py` caller. Additional static reading covered the README, Adversary, InnerLoop, OuterLoop, full PPO driver, attack/environment configurations, and trigger construction. Model utilities and environment wrappers were not exhaustively audited. Full inventory and claim-level coverage are in and .\par
Parent check script SHA-256 values are\par
\section*{Supporting checks}
Native performed the initial independent review and selected source checks. Native challenged the core counterexample and then checked the global reward-domain issue. Their scientific tasks completed and were integrated. Additional read-only initial integration QA, if recorded in the final provenance, is not post-batch acceptance. No person is represented as having reviewed or endorsed these findings.\par
No training, dataset/weight acquisition, author-code execution, paid API, author contact, publication, or external writer was used. An OpenReview attachment was unavailable in the staged acquisition; the embedded appendix was accessible. No accessible Space reference was supplied, and no cloud copy is asserted. Coverage is bounded to the enumerated work and does not certify whole-paper correctness.\par
\end{document}
